\section{Related Work}
\label{sec:related}

\paragraph{Simultaneous speech translation.}
Simultaneous ST couples a translation model with a READ/WRITE policy, from fixed
wait-$k$ schedules \citep{ma2019stacl} to agreement- and attention-based adaptive
policies \citep{liu2020lowlatency,papi2023alignatt}, evaluated by quality jointly
with Average Lagging \citep{ma2019stacl,ma2020simuleval}. Recent work extends SST
to long-form and unbounded streams with computation-aware latency
\citep{ouyang2025infinisst,papi2026doa}. All of these systems are audio-only: at
commit time the model's knowledge of the future is limited to what it can guess
from the source prefix. Our work adds exactly the modality that carries future
information in talk settings.

\paragraph{Visual context for simultaneous \emph{text} MT.}
A line of work initiated by \citet{caglayan2020simultaneous} shows that a paired
image improves low-latency wait-$k$ text MT because ``visual information can
compensate for the missing source context,'' with follow-ups using RL policies
\citep{ive2021exploiting}, multimodal encoders \citep{imankulova2020multimodal},
and supervised visual attention \citep{haralampieva2022supervised}. These studies
establish the \emph{anticipation} value of vision in simultaneous translation, but
in a toy regime: text input (no speech), single caption-style photographs
(Multi30k), and sentence-level decoding. We transfer this insight to its natural
deployment habitat --- long-form talks where the visual channel is a slide stream
that \emph{temporally precedes} the speech discussing it.

\paragraph{Slides and screen text for speech tasks.}
Slide-rich corpora --- SlideSpeech \citep{wang2023slidespeech}, Chinese-LiPS
\citep{zhao2025chineselips}, M3AV \citep{chen2024m3av}, and alignment tooling such
as MaViLS \citep{anderer2024mavils} --- established that projected slides carry
transcription-relevant content. \citet{sinhamahapatra2025slides} demonstrate on
ACL conference talks that slide-derived terminology reduces domain-term WER by
35\% relative, and Slide-EC applies slide terms to post-hoc ASR correction
\citep{piao2025slideec}. All of this work is \emph{offline ASR}: no translation,
no latency dimension, and no policy for deciding when visual evidence should be
trusted. Notably, slide text frequently appears in English even when the talk is
not --- so for X$\rightarrow$En translation (unlike ASR), slides can supply the
\emph{target-side} surface form of upcoming terms, a signal unavailable to
transcription systems by construction.

\paragraph{Multimodal simultaneous ST systems.}
Closest to our setting, OmniFusion \citep{koneru2025omnifusion} fuses a multimodal
foundation model with a translation LLM and evaluates En$\rightarrow$\{De,It\}
SimulST on scientific talks from MCIF \citep{papi2025mcif}. Its latency gain
($\sim$1\,s) is architectural --- removing the ASR$\rightarrow$MT cascade --- and
its visual input is processed \emph{synchronously on the decoding path}: the
authors report that ``image processing introduces additional delay.'' It is
evaluated only against a self-constructed cascade of the same components, assumes
an English source, and neither models the temporal lead of slides nor measures
faithfulness of visual-evidence use. We differ on each axis: X$\rightarrow$En,
an \emph{asynchronous} slide channel off the audio critical path (zero added
per-token latency), an explicit evidence-gating policy, and metrics that separate
vision-induced latency gains from architectural ones.

\paragraph{Audio-visual (lip-based) speech translation.}
MuAViC \citep{anwar2023muavic} and AV-TranSpeech \citep{huang2023avtranspeech} use
lip movements as \emph{phonetic} visual evidence for robust offline ST. Slides
provide \emph{semantic} evidence instead --- readable terms and structure --- which
survives low video quality, requires no face visibility, and is interpretable for
faithfulness auditing; the two channels are complementary.

\paragraph{Retrieval-augmented SST and terminology.}
RASST \citep{luo2026rasst} shows that injecting retrieved glossary terms into SST
improves terminology translation, but assumes the glossary is \emph{prepared in
advance} --- rarely true in open-domain deployment. Our visual channel is the
missing supplier: an on-the-fly glossary harvested from the slide stream,
turning retrieval-augmented SST into a self-provisioning system. Document-level
context for (re-)translation \citep{zhang2021beyond} is similarly complementary:
we focus on context that is external to the speech stream and arrives with
temporal lead.

\paragraph{Positioning.}
No prior work (i) treats vision as a zero-latency \emph{lookahead} channel for
\emph{speech} translation, (ii) evaluates visual evidence under a quality--latency
Pareto rather than offline quality alone, or (iii) measures faithfulness of visual
grounding (visible-but-unspoken hallucination, wrong-slide adoption) in a
streaming regime. These three gaps define our contribution.
